
\begin{frame}[fragile]{Exemplo (Sequência de Fibonacci)}
  \begin{minipage}[t]{0.45\textwidth}
    Fibonacci em Haskell:
    \vspace{0.2cm}
      \begin{code}
f = 0 : g
g = 1 : zipWith (+) f g
  \end{code}
  \end{minipage}
  \pause
  \hfill
  \vrule
  \hfill
  \begin{minipage}[t]{0.50\textwidth}
    Fibonacci em Stream Calculus:
    \vspace{0.2cm}
    \begin{code}
$f$ = $0 : g$
$g$ = $1 : f + g$
    \end{code}
  \end{minipage}
  \vfill
  \pause
  Definimos as operações ($+$), ($-$), ($\cdot$) e ($/$) índice-a-índice, como \var{zipWith} em Haskell.
\end{frame}

\begin{frame}[fragile]{Stream Calculus}
  Definimos operadores ``acumuladores'':
  \vfill
  \pause
  \begin{minipage}[t]{0.45\textwidth}
    Operador $\Sigma$:
    \vspace{0.1cm}
      \[\Sigma a = (0 : a + \Sigma a)\]
  \end{minipage}
  \pause
  \hfill
  \vrule
  \hfill
  \begin{minipage}[t]{0.50\textwidth}
    Operador $\Pi$:
    \vspace{0.1cm}
    \[\Pi a = (1 : a \cdot \Pi a)\]
  \end{minipage}
  \vfill
  \pause
  \begin{minipage}[t]{0.45\textwidth}
    Propriedade Universal do $\Sigma$:
    \vspace{0.1cm}
      \[b = \Sigma a \iff b = (0 : a + b)\]
  \end{minipage}
  \pause
  \hfill
  \vrule
  \hfill
  \begin{minipage}[t]{0.50\textwidth}
    Propriedade Universal do $\Pi$:
    \vspace{0.1cm}
    \[b = \Pi a \iff  b = (1 : a \cdot b)\]
  \end{minipage}
  \vfill
  \pause
  \begin{center}
    Equivalentes à \var{scanl} em Haskell.
  \end{center}
\end{frame}